\scanpage{96}
\uline{Prop} Soit $K$ un corps de nbs fini sur $Q$,
$L$ une extension quadratique de $K$, $p_0$ un idéal
premier de $K$, \deleted{$p,p'$} qui se décompose dans $L$ dans
deux idéaux premiers distincts $p,p'$. Alors il
existe un $x\in L^*$ \deleted{dont le diviseur \unsure{absolu} soit}
\deleted{tel que} dont le diviseur \deleted{soit} $h(p-p')$, \inserted{où $h\in\mathbb N^*$,} et tel que
\deleted{\unsure{Alors}} $x\,\sigma x=1$, où $\sigma$ est l'autom non trivial de $L/K$.
\margin{Plus précisément, à torsion près les diviseurs des éléments de $L$ tels que $N_{L/K}x=1$, sont les diviseurs de $L$ de norme $1$ !}
% slanted

En effet, \deleted{\unsure{comme le groupe des}}\inserted{si $h$ est l'ordre} du groupe de
classes \inserted{de div.} de $L$, il existe un $y$ de diviseur
$h(p-p')$. Alors $y\,\sigma y$ est un élément de $K$ dont le
diviseur est \unsure{évidemt|évidemment} nul, donc c'est une unité $u$.
Donc si $z=y^2$, alors $\deleted{\unsure{x\bar x=}}\ \deleted{\unsure{x\bar x}}\ z\cdot\unsure{{}^{\sigma}z|\bar z}=(y\cdot{}^\sigma y)^2=u^2$,
donc si $x=\frac zu=\frac{y^2}{Ny}$, on aura $\operatorname{div}x=\operatorname{div}z=2\operatorname{div}y$
$=2h(p-p')$, et $x\,{}^\sigma x=(z\,{}^\sigma z)/u^2=1$.

\uline{N.B.} Supposons que $K$ \unsure{ext.|est} tot. réelle \deleted{\unsure{et $L$ une}}
\deleted{\unsure{extension de $Q$,|galoisienne de $Q$,}} \deleted{\unsure{telle que}} ext. quadratique
\deleted{\unsure{de $K$|et $K$}} tot. imaginaire \inserted{de $K$, supposée \uline{galoisienne} sur $Q$,}
. Soit $p$ un nb premier, et
$p_0$ idéal de $K$ au dessus de $p$, supposons que $\unsure{p_0|p}$ il n'y ait
qu'un \uline{seul idéal} \inserted{premier} $p$ \deleted{\unsure{p}} de $L$ sur $K$. Soit $\lambda$ un
\uline{entier} de $L$ tel que $\lambda\bar\lambda=p^i$. Alors $\lambda^2\bar\lambda^2=p^{2i}$ donc
si $\mu=\lambda^2/p^i=\frac{\lambda^2}{N_{L/K}\lambda}$ on aura $\mu\bar\mu=1$, \deleted{et $\mu$ est une unité en}
qui implique que $\mu$ est une unité en $p$, donc ($K$ étant galoisien)
en $\forall$ idéal premier au dessus de $p$ par le \unsure{raisonnt} (l'hyp. sur
$p$ \deleted{\unsure{donc}} \unsure{encore} \deleted{\unsure{avec}} \unsure{vraie} pour les $p'$), donc $\mu$ est une unité (NB. $\lambda$ \unsure{hors} $p$ est une
une racine de l'unité. Donc $\lambda=p^{i/2}u$, $u$ racine
de l'unité !
% archive sheet number: 95

